\chapter{Basis and Characteristic}
\label{ch:basis}

\section*{The Basis Mode}

\begin{versebox}
The Conqueror teaches the teaching ---\\
every state rests on a ground.\\
Trace each thing to what supports it,\\
and the structure will be found.
\end{versebox}

\noindent Every mental state has a parent.\footnote{Pali: \textit{``Tattha katamo padaṭṭhāno hāro? `Dhammaṃ deseti jino'ti, ayaṃ padaṭṭhāno hāro. Kiṃ deseti?''} The mnemonic fragment is ``The Conqueror teaches the teaching'' (\textit{dhammaṃ deseti jino}). The word \textit{padaṭṭhāna} literally means ``standing-place for the feet'' --- a foundation, platform, proximate cause. In Abhidhamma usage it becomes a technical term for the immediate supporting condition of a mental state. \textbf{A note on our translation approach}: the original Pali for this entire section is a systematic catalogue --- ``X has such-and-such a characteristic; its basis is Y'' --- repeated for each of the eighteen root terms and beyond. We have recast this catalogue as narrative prose for readability, grouping related items and using analogies. The full technical listing is preserved in these endnotes for scholars and serious students. The meaning is identical; only the packaging has changed.}

That is the entire insight of the Basis Mode, and it is more powerful than it sounds. Think of it this way. You are trying to understand why you feel a certain way --- angry, anxious, restless, craving something you know is bad for you. Most people stop at the feeling itself: \textit{I'm angry}. But the Basis Mode says: no, go one level deeper. What is the anger standing on? What gave rise to it? Find that, and you have found the lever.

The Buddha's teaching, according to this mode, is essentially a vast map of these parent-child relationships. Every unwholesome state can be traced to its supporting condition. Every wholesome state, too. And the map is precise.

\section*{How a Mind Goes Wrong}

Start with ignorance --- the failure to see things as they really are. What supports it? The distortions: the habit of seeing permanence where there is none, pleasure where there is pain, selfhood where there is no self, beauty where there is foulness. Those four distortions are the ground ignorance stands on.\footnote{Pali: \textit{``Sabbadhammayāthāvaasampaṭivedhalakkhaṇā avijjā, tassā vipallāsā padaṭṭhānaṃ.''} The four distortions (\textit{vipallāsā}) --- perceiving permanence, pleasure, self, and beauty where none exist --- operate at three levels: perception (\textit{saññā}), thought (\textit{citta}), and view (\textit{diṭṭhi}). See AN 4.49 for the canonical treatment.}

Now take craving --- the mind's compulsive clinging. What feeds it? Whatever is pleasant, whatever is agreeable. That is the ground craving stands on. You don't crave what repels you; craving needs something sweet to latch onto.\footnote{Pali: \textit{``Ajjhosānalakkhaṇā taṇhā, tassā piyarūpaṃ sātarūpaṃ padaṭṭhānaṃ.''} The phrase \textit{piyarūpaṃ sātarūpaṃ} --- ``what has the form of the dear, what has the form of the pleasant'' --- recurs throughout the Mahānidāna Sutta (DN 15) as the trigger for craving. The basis of craving is not the pleasant object itself but one's \textit{perception} of it as pleasant.}

Greed stands on the habit of taking what isn't yours. The fixation on physical beauty stands on the failure to guard your senses --- you let your eyes linger, and the obsession has its foothold. The illusion that things are pleasurable stands on the gratification you get from sensory contact. The illusion that things are permanent stands on the continuity of consciousness itself --- because consciousness seems to flow unbroken, you mistake the river for something solid. And the illusion of a self? That stands on the whole mental apparatus --- the aggregate of ``naming'' that creates the narrative of ``me.''\footnote{Pali: \textit{``Patthanalakkhaṇo lobho, tassa adinnādānaṃ padaṭṭhānaṃ. Vaṇṇasaṇṭhānabyañjanaggahaṇalakkhaṇā subhasaññā, tassā indriyāsaṃvaro padaṭṭhānaṃ. Sāsavaphassaupagamanalakkhaṇā sukhasaññā, tassā assādo padaṭṭhānaṃ. Saṅkhatalakkhaṇānaṃ dhammānaṃ asamanupassanalakkhaṇā niccasaññā, tassā viññāṇaṃ padaṭṭhānaṃ. Aniccasaññādukkhasaññāasamanupassanalakkhaṇā attasaññā, tassā nāmakāyo padaṭṭhānaṃ.''} Here is the full technical mapping for the unwholesome side: (1) greed (\textit{lobha}, characteristic: wanting) → basis: taking what is not given; (2) beauty-perception (\textit{subhasaññā}, characteristic: grasping at form and shape) → basis: failure to restrain the senses; (3) pleasure-perception (\textit{sukhasaññā}, characteristic: approaching corrupt contact) → basis: gratification; (4) permanence-perception (\textit{niccasaññā}, characteristic: not seeing the conditioned marks) → basis: consciousness; (5) self-perception (\textit{attasaññā}, characteristic: not seeing impermanence and suffering) → basis: the mental body (\textit{nāmakāya}).}

Each of these is like a weed. Pull the weed and it grows back. But dig out its root --- its basis --- and the ground is clear.

\section*{How a Mind Goes Right}

The wholesome side is a mirror image. True knowledge --- the capacity to penetrate all things --- stands on everything that can be known. It needs the whole world as its classroom. Calm --- the ability to gather a scattered mind --- stands on contemplating the unattractive, because nothing quiets a restless, craving mind like seeing things without their glamour.\footnote{Pali: \textit{``Sabbadhammasampaṭivedhalakkhaṇā vijjā, tassā sabbaṃ neyyaṃ padaṭṭhānaṃ. Cittavikkhepapaṭisaṃharaṇalakkhaṇo samatho, tassa asubhā padaṭṭhānaṃ.''} True knowledge (\textit{vijjā}) has quite literally ``everything that can be known'' (\textit{sabbaṃ neyyaṃ}) as its basis. Calm (\textit{samatha}) is grounded in the contemplation of the unattractive (\textit{asubha}) --- a meditation practice in which the meditator reflects on the body's lack of inherent beauty.}

Non-greed stands on the practice of not taking what isn't given. Non-hatred stands on the practice of not killing. Non-delusion stands on right practice itself. And the corrective perceptions --- seeing things as unattractive, as painful, as impermanent, as not-self --- each stands on its own specific ground: disenchantment, feeling, arising-and-passing-away, and the perception of things as they actually are.\footnote{Pali: \textit{``Icchāvacarapaṭisaṃharaṇalakkhaṇo alobho, tassa adinnādānā veramaṇī padaṭṭhānaṃ. Abyāpajjalakkhaṇo adoso, tassa pāṇātipātā veramaṇī padaṭṭhānaṃ. Vatthuavippaṭipattilakkhaṇo amoho, tassa sammāpaṭipatti padaṭṭhānaṃ. Vinīlakavipubbakagahaṇalakkhaṇā asubhasaññā, tassā nibbidā padaṭṭhānaṃ. Sāsavaphassaparijānanalakkhaṇā dukkhasaññā, tassā vedanā padaṭṭhānaṃ. Saṅkhatalakkhaṇānaṃ dhammānaṃ samanupassanalakkhaṇā aniccasaññā, tassā uppādavayā padaṭṭhānaṃ. Sabbadhammaabhinivesalakkhaṇā anattasaññā, tassā dhammasaññā padaṭṭhānaṃ.''} The full technical mapping for the wholesome side. Notice the elegant symmetry: each unwholesome perception is mirrored by a wholesome perception with the opposite characteristic, and the basis of the wholesome perception is precisely what uproots the unwholesome one. For example, the \textit{permanence}-perception (unwholesome) stands on not-seeing arising and passing away; the \textit{impermanence}-perception (wholesome) stands on seeing arising and passing away. Same ground, opposite directions.}

Notice the pattern. Every unwholesome state has a mirror on the wholesome side, and the basis of the cure is always the direct opposite of the basis of the disease. It is not random. It is not mystical. It is a systematic map of the mind's architecture, and once you see it, you can navigate.

\section*{Three Kinds of Craving, Three Kinds of Fuel}

The map goes further. Sensual craving --- the craving for sights, sounds, smells, tastes, and touches --- stands on the five strands of sense-pleasure. Craving for form --- the subtler attachment to the beauty of the meditative realms --- stands on the five physical sense faculties. And craving for existence itself --- the deepest, most tenacious clinging --- stands on the mind, the sixth sense-base.\footnote{Pali: \textit{``Pañca kāmaguṇā kāmarāgassa padaṭṭhānaṃ, pañcindriyāni rūpīni rūparāgassa padaṭṭhānaṃ, chaṭṭhāyatanaṃ bhavarāgassa padaṭṭhānaṃ.''} Three cravings, three domains: sense-desire (\textit{kāma}), form (\textit{rūpa}), and existence (\textit{bhava}). Each has its own specific fuel.}

Each craving has its own fuel supply. Cut the fuel, and the fire goes out --- which is, of course, exactly what the word ``nirvana'' means.\footnote{The etymological connection between \textit{nibbāna} (nirvana) and the extinguishing of fire is well established. The metaphor is central to early Buddhist thought: defilements are the fire; their conditions are the fuel (\textit{upādāna}, which also means ``clinging''); and liberation is what happens when the fuel is removed. See Chapter~\ref{ch:fitness} on the fire analogy for craving.}

\section*{Faith, Energy, Mindfulness, Concentration, Wisdom}

What about the five spiritual faculties --- the core capacities that drive practice? Each rests on its own specific training ground.

Faith --- the ability to commit --- stands on verified confidence: you believe because you have \textit{tested} and found reliable. And confidence --- that inner clarity, like clear water --- stands on faith. The two support each other like two hands washing.\footnote{Pali: \textit{``Okappanalakkhaṇā saddhā adhimuttipaccupaṭṭhānā ca, anāvilalakkhaṇo pasādo sampasīdanapaccupaṭṭhāno ca. Abhipatthiyanalakkhaṇā saddhā, tassā aveccapasādo padaṭṭhānaṃ. Anāvilalakkhaṇo pasādo, tassa saddhā padaṭṭhānaṃ.''} The Netti distinguishes two aspects of faith: the aspiring aspect (\textit{abhipatthiyana}) and the clarifying aspect (\textit{pasāda}, ``serene confidence,'' literally ``clarity, non-turbidity''). They are mutually supportive: trust leads to clarity, and clarity reinforces trust.}

Energy stands on the four right strivings --- the practice of actually showing up and doing the work. Mindfulness stands on the four foundations of mindfulness --- the practice of paying close attention to body, feeling, mind, and mental qualities. Concentration stands on the absorptions --- the deep states of meditative unification. And wisdom stands on the truths --- the Four Noble Truths, the bedrock of the entire teaching.\footnote{Pali: \textit{``Ārambhalakkhaṇaṃ vīriyaṃ, tassa sammappadhānaṃ padaṭṭhānaṃ. Apilāpanalakkhaṇā sati, tassā satipaṭṭhānaṃ padaṭṭhānaṃ. Ekaggalakkhaṇo samādhi, tassa jhānāni padaṭṭhānaṃ. Pajānanalakkhaṇā paññā, tassā saccāni padaṭṭhānaṃ.''} The Pali characterization of mindfulness is especially vivid: \textit{apilāpana}, ``non-floating.'' Mindfulness is defined as the quality of not drifting --- it stays, it sinks in, it does not bob on the surface.}

\section*{The Domino Chain}

Now the Basis Mode reveals its most dramatic insight. There is another approach, the text says --- and then it lays out the entire twelve-link chain of dependent origination, reframed as a chain of bases.\footnote{Pali: \textit{``Aparo nayo, assādamanasikāralakkhaṇo ayonisomanasikāro, tassa avijjā padaṭṭhānaṃ. Saccasammohanalakkhaṇā avijjā, sā saṅkhārānaṃ padaṭṭhānaṃ.''} This ``other approach'' (\textit{aparo nayo}) reframes the twelve links of dependent origination (\textit{paṭiccasamuppāda}) as a chain of \textit{padaṭṭhāna}. It's a creative move: dependent origination is not just a causal chain but a chain of standing-places --- each link is the platform the next one stands on.}

Think of dominoes. Unwise attention --- the habit of focusing on what's gratifying --- is the platform ignorance stands on. Ignorance is the platform that mental formations stand on. Formations are the platform consciousness stands on. And on it goes: consciousness supports name-and-form, which supports the six senses, which support contact, which supports feeling, which supports craving, which supports clinging, which supports existence, which supports birth.

And then the sequence accelerates into darkness. Birth is the platform for aging. Aging is the platform for death. Death is the platform for sorrow. Sorrow produces agitation, which is the platform for lamentation. Lamentation produces wailing, which is the platform for pain. Pain crushes the body, which is the platform for grief. Grief crushes the mind, which is the platform for despair. And despair --- this is the devastating part --- despair produces a kind of burning, and that burning is the platform for existence.\footnote{Pali: \textit{``Khandhapātubhavanalakkhaṇā jāti, sā jarāya padaṭṭhānaṃ \ldots Odahanakārako upāyāso, so bhavassa padaṭṭhānaṃ.''} The Pali gives a vivid characteristic for each link. Aging is ``the ripening of acquisitions'' (\textit{upadhiparipāka}); death is ``the cutting-off of the life-faculty'' (\textit{jīvitindriyupaccheda}); sorrow ``produces agitation'' (\textit{ussukkakāraka}); lamentation ``produces wailing'' (\textit{lālappakāraka}); despair ``produces burning'' (\textit{odahanakāraka}). And then the chain circles back: despair feeds right back into existence.}

The circle closes. Despair feeds back into existence. The last domino tips the first. When all these constituents arise together, that is existence --- and existence is the basis of the entire round of rebirth.\footnote{Pali: \textit{``Imāni bhavaṅgāni yadā samaggāni nibbattāni bhavanti so bhavo, taṃ saṃsārassa padaṭṭhānaṃ.''}}

But there is one domino that tips \textit{outward} instead of around. The path --- which has the characteristic of leading out --- is the basis of cessation.\footnote{Pali: \textit{``Niyyānikalakkhaṇo maggo, so nirodhassa padaṭṭhānaṃ.''} One sentence to close the circle of suffering; one sentence to break it. The path (\textit{magga}) is the only link in the chain that leads \textit{out} rather than \textit{around}.}

One sentence to describe the trap. One sentence to describe the exit.

\section*{The Upward Spiral}

And then the Netti does something beautiful. Having shown the downward spiral, it traces an upward one --- a chain of conditions that leads not to suffering but to freedom. And this one reads like music.\footnote{Pali: \textit{``Titthaññutā pītaññutāya padaṭṭhānaṃ \ldots vimutti vimuttiñāṇadassanassa padaṭṭhānaṃ.''} This chain appears in numerous Nikāya passages (e.g., AN 10.1, AN 11.1) and represents the natural unfolding of the path when its conditions are met. The Netti's innovation is to frame it as a \textit{padaṭṭhāna} sequence --- each quality is the standing-place for the next.}

Knowing where you stand is the basis of knowing what you've taken in. Knowing what you've taken in is the basis of knowing your capacity. Knowing your capacity is the basis of knowing yourself. Knowing yourself is the basis of having made merit in the past. Past merit is the basis of living in the right place. Living in the right place is the basis of finding good people. And finding good people is the basis of pointing yourself in the right direction.

Then the chain lifts off. Pointing yourself in the right direction is the basis of virtue. Virtue is the basis of a clear conscience. A clear conscience is the basis of delight. Delight is the basis of rapture. Rapture is the basis of tranquility. Tranquility is the basis of happiness. Happiness is the basis of concentration. Concentration is the basis of seeing things as they really are. Seeing things as they really are is the basis of disenchantment. Disenchantment is the basis of dispassion. Dispassion is the basis of liberation. And liberation is the basis of \textit{knowing you are free}.

Every supporting condition, every cause --- all of it is a basis. That is why the Venerable Mahākaccāyana said: ``The Conqueror teaches the teaching.''\footnote{Pali: \textit{``Evaṃ yo koci upanissayo yo koci paccayo, sabbo so padaṭṭhānaṃ. Tenāha āyasmā mahākaccāyano `dhammaṃ deseti jino'ti.''}}

\ornament

\section*{The Characteristic Mode}

\begin{versebox}
When one thing has been spoken of,\\
and others share its mark ---\\
then all the rest are spoken too:\\
one light can fill the dark.
\end{versebox}

\noindent Here is a principle so simple it sounds like a trick --- but it is one of the most powerful reading tools in the entire Netti.\footnote{Pali: \textit{``Tattha katamo lakkhaṇo hāro? `Vuttamhi ekadhamme'ti, ayaṃ lakkhaṇo hāro. Kiṃ lakkhayati? Ye dhammā ekalakkhaṇā, tesaṃ dhammānaṃ ekasmiṃ dhamme vutte avasiṭṭhā dhammā vuttā bhavanti.''} The mnemonic fragment is ``When one thing has been spoken of'' (\textit{vuttamhi ekadhamme}). The principle: shared characteristic = shared applicability.}

If two things share a defining characteristic, then what is said about one has been said about the other.

That is it. That is the whole mode. But watch what it does.

\section*{The Murderer Test}

The Buddha once described the eye in devastating terms: ``The eye, monks, is unstable, fleeting, limited, fragile, alien, painful, ruinous, shaky, a furnace, a conditioned thing, a \textit{murderer}, surrounded by enemies.'' Twelve adjectives, and every one of them is true of the eye.\footnote{Pali: \textit{``Cakkhuṃ, bhikkhave, anavaṭṭhitaṃ ittaraṃ parittaṃ pabhaṅgu parato dukkhaṃ byasanaṃ calanaṃ kukkuḷaṃ saṅkhāraṃ vadhakaṃ amittamajjhe.''} Twelve adjectives. The Netti zeroes in on the key one: \textit{vadhaka}, ``murderous.'' This is the shared characteristic that links all six internal sense-bases.}

Now: the ear, the nose, the tongue, the body, and the mind all share the defining characteristic of the eye --- they are all ``murderous'' in the same way. So when the Buddha described the eye, he was \textit{simultaneously} describing all six senses. He did not need to say the same thing six times. One example was enough.

Think of it like discovering that all metals conduct electricity. Once you know that copper conducts, you don't need to test every piece of iron, gold, and silver separately. The shared property does the work for you.

\section*{Five Murderers}

The same principle applies to the five aggregates. The Buddha told the monk Rādha: ``Be unconcerned with past form. Do not delight in future form. Practice for the disenchantment, dispassion, and cessation of present form.'' He said it about form. But the Discourse of Yamaka's Advice compares all five aggregates to five murderers pursuing a man --- they all share the ``murderous'' characteristic. So when the Buddha spoke about form, he simultaneously spoke about feeling, perception, mental formations, and consciousness. One aggregate, five instructions.\footnote{Pali: \textit{``Atīte, rādha, rūpe anapekkho hohi \ldots Sabbe hi pañcakkhandhā yamakovādasutte vadhakaṭṭhena ekalakkhaṇā vuttā.''} The Yamaka Sutta (SN 22.85) uses the famous simile of five murderers, each representing one aggregate. Since all five share the same defining mark, what is said about one applies to all.}

And the Dhammapada verse:

\begin{versebox}
Those who have rightly undertaken\\
constant mindfulness of the body ---\\
they do not practice what should not be done;\\
in what should be done, they persevere.
\end{versebox}

When mindfulness of the body has been spoken of, mindfulness of feelings, mind, and mental qualities have all been spoken of too --- because all four foundations of mindfulness share a single characteristic.\footnote{Pali: \textit{``Iti kāyagatāya satiyā vuttāya vuttā bhavanti vedanāgatā sati cittagatā dhammagatā ca.''} The verse is from Dhp 293.} Likewise, when the Buddha said ``whatever is seen, heard, or sensed'' --- ``cognized'' has been said too. The list didn't need a fourth item; the shared characteristic implied it.\footnote{Pali: \textit{``Tathā yaṃ kiñci diṭṭhaṃ vā sutaṃ vā mutaṃ vā'ti vutte vuttaṃ bhavati viññātaṃ.''} The standard formula \textit{diṭṭha-suta-muta-viññāta} (seen-heard-sensed-cognized) treats the first three as exemplary; the fourth is always implied.}

\section*{A Single Sentence, Four Practices}

Now the Netti pulls off a stunning piece of close reading. Take the most famous meditation instruction in the entire Canon: ``Therefore, monk, dwell contemplating the body in the body --- ardent, clearly comprehending, mindful, having removed longing and grief regarding the world.''

Four adjectives. Four entire training systems, packed into a single sentence:\footnote{Pali: \textit{``Ātāpī'ti vīriyindriyaṃ, `sampajāno'ti paññindriyaṃ, `satimā'ti satindriyaṃ, `vineyya loke abhijjhādomanassan'ti samādhindriyaṃ.''}}

``Ardent'' = the faculty of energy.\\
``Clearly comprehending'' = the faculty of wisdom.\\
``Mindful'' = the faculty of mindfulness.\\
``Having removed longing and grief'' = the faculty of concentration.

When you practice body-contemplation with these four qualities, all four foundations of mindfulness reach their fulfillment simultaneously. Why? Because the four faculties share a single characteristic. They are not four separate things bolted together. They are one practice seen from four angles.\footnote{Pali: \textit{``Evaṃ kāye kāyānupassino viharato cattāro satipaṭṭhānā bhāvanāpāripūriṃ gacchanti. Kena kāraṇena, ekalakkhaṇattā catunnaṃ indriyānaṃ.''} This is the Characteristic Mode at its most practical: it shows that the four foundations of mindfulness are not four separate practices to be done sequentially, but four aspects of a single practice that arise together.}

\section*{The Domino Effect of Awakening}

The cascade goes further --- much further. When the four foundations of mindfulness are developed, the four right strivings reach fulfillment. When the right strivings are developed, the four bases of power reach fulfillment. When those are developed, the five faculties reach fulfillment. Then the five powers. Then the seven factors of awakening. Then the Noble Eightfold Path.\footnote{Pali: \textit{``Catūsu satipaṭṭhānesu bhāviyamānesu cattāro sammappadhānā bhāvanāpāripūriṃ gacchanti \ldots ariyo aṭṭhaṅgiko maggo bhāvanāpāripūriṃ gacchati.''} The thirty-seven qualities conducive to awakening (\textit{bodhipakkhiyā dhammā}): 4 + 4 + 4 + 5 + 5 + 7 + 8 = 37. All share the single characteristic of ``leading out'' (\textit{niyyānika}).}

Thirty-seven qualities. One shared characteristic: they all lead out. And because they share that one characteristic, developing any one of them is developing all of them. You do not need to practice thirty-seven separate things. You practice \textit{one} thing that naturally manifests in thirty-seven ways.\footnote{Pali: \textit{``Sabbe hi bodhaṅgamā bodhipakkhiyā neyyānikalakkhaṇena ekalakkhaṇā, te ekalakkhaṇattā bhāvanāpāripūriṃ gacchanti.''} This is arguably the Characteristic Mode's single most important practical claim: the entire path is a unity. Like a prism that separates white light into a spectrum, the thirty-seven factors are not different lights but one light viewed through different facets.}

\section*{The Flip Side: Abandonment}

The same domino effect works for getting rid of what harms you. When the foundations of mindfulness are developed, the distortions fall away. The nutriments of suffering are understood. You become free of clinging, unyoked from bondage, released from the knots that tie you up, free of the corruptions that leak into everything, crossed over the floods that drown you, free of the darts that pierce you, no longer taking wrong courses.\footnote{Pali: \textit{``Catūsu satipaṭṭhānesu bhāviyamānesu vipallāsā pahīyanti, āhārā cassa pariññaṃ gacchanti, upādānehi anupādāno bhavati, yogehi ca visaṃyutto bhavati, ganthehi ca vippayutto bhavati, āsavehi ca anāsavo bhavati, oghehi ca nitthiṇṇo bhavati, sallehi ca visallo bhavati, viññāṇaṭṭhitiyo cassa pariññaṃ gacchanti, agatigamanehi na agatiṃ gacchati.''} The Pali lists ten categories of defilements, each with its own standard metaphor: distortions (\textit{vipallāsa}), nutriments (\textit{āhāra}), clingings (\textit{upādāna}), yokes (\textit{yoga}), knots (\textit{gantha}), corruptions (\textit{āsava}), floods (\textit{ogha}), darts (\textit{salla}), stations of consciousness (\textit{viññāṇaṭṭhiti}), and wrong courses (\textit{agati}). Ten different ways of naming the same bondage; one practice dissolves them all.}

All unwholesome states share a single characteristic: they are what gets abandoned when mindfulness is developed. So they all go --- together, in one stroke, like ten different names for the same prisoner being released.

\section*{The Hidden Depths of Neutral Feeling}

One last example, and it is the most surprising. Wherever pleasant feeling has been taught, the Netti says, the truth of the origin of suffering has been taught too --- because pleasant feeling is what craving latches onto. And wherever painful feeling has been taught, the truth of suffering itself has been taught.\footnote{Pali: \textit{``Yattha vā pana sukhā vedanā desitā, desitaṃ tattha sukhindriyaṃ somanassindriyaṃ dukkhasamudayo ca ariyasaccaṃ. Yattha vā pana dukkhā vedanā desitā, desitaṃ tattha dukkhindriyaṃ domanassindriyaṃ dukkhañca ariyasaccaṃ.''}}

But here is the stunning claim: wherever neither-painful-nor-pleasant feeling has been taught, \textit{the entire chain of dependent origination} has been taught.\footnote{Pali: \textit{``Yattha vā pana adukkhamasukhā vedanā desitā, desitaṃ tattha upekkhindriyaṃ sabbo ca paṭiccasamuppādo. Kena kāraṇena, adukkhamasukhāya hi vedanāya avijjā anuseti.''} Neutral feeling contains the entire mechanism of suffering. Why? Because ignorance lies latent (\textit{anuseti}) specifically in neutral feeling --- it is the feeling we do not notice, and therefore the one where ignorance operates most freely.}

Why? Because ignorance lies latent in neutral feeling. It is the feeling you barely notice --- the background hum of ``nothing special'' --- and precisely because you don't notice it, ignorance operates there most freely. And where ignorance lurks, the whole chain can unspool: formations, consciousness, name-and-form, the six senses, contact, feeling, craving, clinging, existence, birth, aging and death, and the entire mass of suffering.

The most innocuous feeling in the world contains, in seed form, the entire mechanism of suffering. And the remedy? The noble qualities that are free from lust, hatred, and delusion.\footnote{Pali: \textit{``So ca sarāgasadosasamohasaṃkilesapakkhena hātabbo, vītarāgavītadosavītamohaariyadhammehi hātabbo.''} The full dependent origination formula is deployed here to make a single point: even the most ``neutral'' experience is not spiritually neutral. The Characteristic Mode reveals the depths beneath apparently simple surfaces.}

\ornament

In this way, whatever things share a single characteristic --- whether in function, in defining mark, in commonality, or in their pattern of passing away and arising --- when one of those things has been spoken of, all the rest have been spoken of too. That is why the Venerable Mahākaccāyana said: ``When one thing has been spoken of.''\footnote{Pali: \textit{``Evaṃ ye dhammā ekalakkhaṇā kiccato ca lakkhaṇato ca sāmaññato ca cutūpapātato ca, tesaṃ dhammānaṃ ekasmiṃ dhamme vutte avasiṭṭhā dhammā vuttā bhavanti.''} Four types of shared characteristic: function (\textit{kicca}), defining mark (\textit{lakkhaṇa}), commonality (\textit{sāmañña}), and pattern of death-and-rebirth (\textit{cutūpapāta}). The Characteristic Mode is the most powerful lens in the Netti's toolkit: it turns every specific statement into a universal one.}
